\chapter{Methods\label{cha:chapter5}}

\section{Research Design and Study Area\label{sec:research_methods}}

\subsection{Overall Research Design}

This study followed a retrospective, predictive research design that linked satellite-based SIF enhancement with an agricultural yield application. The work was organized into two connected analytical stages rather than a single end-to-end model. In the first stage, machine-learning models learned relationships between OCO-2 SIF observations and spatially complete optical, biophysical, radiation, seasonal, and crop-type predictors. In the second stage, the enhanced SIF produced by the best-performing spatial model was restricted to winter-wheat pixels, summarized by month and administrative region, and evaluated as a predictor of reported winter-wheat yield. The first stage therefore addressed the reconstruction and spatial enhancement of a sparsely observed environmental variable, whereas the second stage tested whether the enhanced signal provided useful information for a concrete agricultural task. Both stages used historical observations and held-out data, so the reported results describe predictive association and generalization rather than a causal effect of any individual predictor. The experimental units changed with the question being tested, ranging from individual OCO-2 footprints and fixed spatial cells to NUTS-3 region-year records. This nested design made it possible to investigate model accuracy, input representation, spatial aggregation, agro-climatic variation, and crop-specific utility within one coherent workflow.

The SIF enhancement problem was expressed generally as a supervised mapping from a set of spatial predictors to a SIF value defined over a particular spatial support. Let $u$ denote that support, $t$ the observation time, $\mathbf{X}_{u,t}$ the associated predictor data, and $S_{u,t}$ the corresponding OCO-2-derived target. The general relationship estimated by each model family was

\begin{equation}
    \widehat{S}_{u,t} = f_{\theta}\!\left(\mathbf{X}_{u,t}\right),
    \label{eq:general_sif_mapping}
\end{equation}

where $f_{\theta}$ represents a fitted tree-based, multilayer-perceptron, or convolutional model with parameters $\theta$. For polygon-mean models, $\mathbf{X}_{u,t}$ was a vector of footprint-level summary statistics and the output was one scalar prediction. For convolutional models, $\mathbf{X}_{u,t}$ was a multi-channel raster chip and the network produced a spatial SIF field before reducing it over one or more footprint masks. For the spatial-aggregate model, the target represented the mean of several same-date OCO-2 footprints assigned to one fixed 4 km cell. The definition of $u$ was therefore not interchangeable across experiments, and all reported comparisons retained the target support as part of the model description. This support-aware formulation was necessary because a low error for a noise-reduced 4 km target does not imply the same performance at the native footprint scale. It also provided the basis for the later external experiment in which a model trained with 4 km supervision was evaluated against individual OCO-2 footprint means on previously unseen tiles.

The SIF experiments varied both model architecture and the representation supplied to the model. The tabular branch represented every footprint by polygon means and fitted Random Forest and XGBoost regressors. A separate area-to-point branch used a multilayer perceptron to relate pixel-level Sentinel-2 information to footprint-scale SIF supervision. Single-footprint convolutional experiments placed one target footprint in each image chip, while multi-footprint experiments placed several same-date footprints in one chip and computed a separate mask-averaged prediction for each footprint. The principal spatial-aggregate experiment instead combined at least five OCO-2 observations within a fixed 4 km grid cell and trained a U-Net-like convolutional model with an aggregate footprint-weight map. MODIS-based and Sentinel-2-based configurations were both retained because they tested the consequences of predictor resolution and spatial context as well as network design. Consequently, the broad comparison between all configurations is an empirical comparison of complete modelling systems rather than a perfectly controlled architecture-only ablation. More direct claims about architecture were based on configurations that shared similar targets and evaluation supports, while broader differences were interpreted jointly with their input representation and data preparation.

The downstream yield analysis used the same general predictive logic at a coarser administrative and annual support. For NUTS-3 region $r$ and year $y$, winter-wheat yield was represented as

\begin{equation}
    \widehat{Y}_{r,y} = g_{\phi}\!\left(S_{r,y,3},S_{r,y,4},S_{r,y,5},S_{r,y,6},S_{r,y,7}\right),
    \label{eq:yield_design}
\end{equation}

where $S_{r,y,m}$ denotes a monthly SIF predictor for March through July and $g_{\phi}$ denotes the fitted yield model. Separate versions of the predictor vector were produced from crop-pure enhanced SIF and from raw, mixed-pixel OCO-2 observations. This design held the yield target, region-year records, month set, and Random Forest configuration constant while changing the SIF representation. The crop-pure predictors were computed only from pixels classified as winter wheat, whereas the raw baselines summarized accepted OCO-2 footprints with either no crop-purity requirement or a minimum winter-wheat coverage threshold. The years 2019--2022 formed the model-development period and 2024 was reserved as a temporal test year. The comparison therefore measured whether crop-targeted enhancement reduced out-of-year yield-prediction error relative to spatially mixed satellite observations. It did not assume that the 20 m outputs were directly observed 20 m SIF; rather, they were model-derived spatial estimates constrained during SIF training by coarser OCO-2 supervision.

\subsection{Geographical Scope}

The primary SIF study domain was Germany, represented by a set of Sentinel-2 Military Grid Reference System (MGRS) tiles rather than by a single wall-to-wall national raster. The nine tiles used for the main Sentinel-2 model-development data were \texttt{T32ULA}, \texttt{T32ULB}, \texttt{T32UMD}, \texttt{T32UNA}, \texttt{T32UNC}, \texttt{T32UNE}, \texttt{T32UPC}, \texttt{T32UPE}, and \texttt{T32UPU}. These tiles sampled agricultural landscapes across the north, centre, and south of Germany and included the major crop-producing states used throughout the study. Tile boundaries were operational units for organizing large raster products, assigning SIF observations, constructing chips, and avoiding accidental mixing of incompatible Sentinel-2 product grids. The spatial distribution of the retained SIF observations and their relationship to the selected MGRS boundaries are shown in Figure~\ref{fig:sif_centroids_mgrs_boundaries}. The design does not claim uninterrupted 20 m coverage of every location in Germany for every experiment, because model-ready observations were also constrained by OCO-2 tracks, product availability, quality filters, and the selected months. Germany-wide conclusions therefore refer to the diversity represented by the selected tiles and observations rather than to an exhaustive census of all German pixels. This distinction was retained when interpreting spatial residuals and subgroup results.

Two additional MGRS tiles, \texttt{T32UQV} and \texttt{T32UMC}, were used for external native-footprint evaluation of the principal 4 km spatial-aggregate model. Neither tile contributed samples to the fitting, validation, calibration, or internal testing of that model. Accepted footprints from these tiles were centred within 4 km by 4 km predictor windows at 20 m resolution, and the resulting SIF map was averaged using the fractional footprint mask. To prevent truncated predictor windows, only observations inside an inner 99.8 km square of each 109.8 km Sentinel-2 tile were eligible for sampling. This removed a 5 km border on every side and provided more margin than the 2 km distance from a centred footprint to the edge of its input window. The external experiment consequently tested both geographic transfer to unseen tiles and transfer from aggregated training support back to approximately native OCO-2 footprint support. It was kept separate from the internal test set because these two evaluations answer different questions. Internal testing quantified performance on held-out acquisition components within the main study tiles, whereas external testing examined whether the fitted relationship remained useful in new spatial domains.

The winter-wheat yield application focused on Bavaria because annual NUTS-3 yield statistics, annual crop-type rasters, and the required monthly Sentinel-2 composites could be aligned for this state. Five MGRS tiles were used for the expanded Bavarian domain: \texttt{T32UNA}, \texttt{T32UPU}, \texttt{T32UQV}, \texttt{T32UPA}, and \texttt{T32UPV}. A Bavarian NUTS-3 region was retained when at least 50\% of its area fell within one of these selected Sentinel-2 tile extents. If a region satisfied the threshold for more than one tile, it was assigned once using a deterministic preference for previously processed tiles and then the largest overlap percentage. This produced non-duplicated administrative units while retaining substantially more of Bavaria than the earlier 80\% overlap criterion. The selected administrative regions, MGRS extents, and associated agro-climatic context are summarized spatially in Figure~\ref{fig:hzs_mgrs_nuts3}. Yield prediction was performed at the NUTS-3 region-year level because this was the common support of the official response data. Pixel-level SIF predictions were therefore treated as intermediate information that had to be restricted to winter wheat and aggregated before being related to yield.

Agro-climatic variation was represented independently of the administrative and raster tiling systems. Each SIF observation inherited one or more plant hardiness zone labels from the spatial zone layer, and model performance was later summarized for the individual zones 7a, 7b, 8a, 8b, and 9a. Mixed labels could occur when an aggregate support intersected more than one zone, but the principal zone comparison reported groups with a single zone label to keep interpretation clear. MGRS tiles, federal states, NUTS regions, and hardiness zones therefore served different purposes and were not treated as interchangeable regional definitions. MGRS tiles defined raster storage and major experimental partitions, federal states supplied broad geographic context, NUTS-3 regions defined the yield response support, and hardiness zones described agro-climatic conditions. Keeping these roles separate reduced the risk of attributing a tile-specific pattern directly to climate or an administrative effect. It also allowed the same fitted SIF model to be evaluated by geography, season, measurement mode, and agro-climatic zone without changing its prediction target. The study area was consequently a hierarchy of overlapping analytical geographies rather than one fixed boundary used for every task.

\subsection{Temporal Scope and Spatial Supports}

The main SIF enhancement datasets covered 2019--2024 and focused on February through July, when the principal growth stages of German crops occur. OCO-2 observations remained tied to individual acquisition dates and were therefore irregular in time. Sentinel-2 predictors were monthly composites, GLASS vegetation products were generally eight-day composites, and VIIRS PAR represented a daily mean. Dates were matched according to the meaning of each product rather than with one common nearest-date rule. Sentinel-2 records were linked by month and product interval, eight-day products were linked to the interval containing the SIF date, and VIIRS PAR was matched by exact date. OCO-2 records just outside a Sentinel-2 product interval were retained with a \texttt{date\_align} flag for later evaluation. The yield analysis used March through July in 2019--2022 and 2024, while 2023 was excluded because the necessary Sentinel-2 coverage was incomplete. Years 2019--2022 were used for model development and 2024 was held out for testing, including when missing-value replacement statistics were calculated. Keeping the months as separate predictors preserved seasonal changes in winter-wheat growth and senescence.

The workflow also crossed several spatial supports that were kept distinct during modelling and evaluation. Sentinel-2 bands at 10 m and 20 m were aligned to a common 20 m grid, while the 10 m crop map was converted to fractional crop information on that grid. GLASS FAPAR had an actual pixel spacing of approximately 231.7 m, and VIIRS PAR had a spacing of approximately 927 m before both were aligned with the Sentinel-2 predictors. OCO-2 observations were irregular polygons with processed areas typically close to 1 km$^2$. Single-footprint and multi-footprint experiments preserved these individual footprint targets, although multi-footprint chips included several separate masks. The spatial-aggregate experiment instead averaged at least five same-date observations within a fixed 4 km cell represented by 200 by 200 predictor pixels. External inference then tested the 4 km-trained model by averaging its 20 m predictions over individual OCO-2 footprints in unseen tiles. For the yield analysis, predictions over winter-wheat pixels were aggregated again to one monthly value for each NUTS-3 region and year. Model results were always compared at the spatial support of the available reference value, so the 20 m output maps were not treated as directly observed 20 m SIF.

\section{Datasets\label{sec:datasets}}

\subsection{OCO-2 SIF Observations}

OCO-2 observations provided the response variable used throughout the SIF enhancement experiments. The source was the Level 2 bias-corrected solar-induced fluorescence product generated with the IMAP-DOAS retrieval and distributed as daily files by the NASA Goddard Earth Sciences Data and Information Services Center. Retrospective product releases V11r and V11.2r were represented in the local archive used for this study \cite{oco2_sif_v11r,oco2_sif_v112r}. The files contained daily-corrected SIF retrievals at 740 nm, 757 nm, and 771 nm, together with retrieval-uncertainty fields for the oxygen absorption bands. Each sounding included its acquisition time, central latitude and longitude, and four latitude--longitude corners describing the elongated ground footprint. Quality flags, measurement mode, viewing geometry, illumination geometry, source-file identifiers, and selected meteorological fields were also available. Measurement mode distinguished nadir and glint observations, while the four corners retained the actual spatial support of each retrieval rather than reducing it to a point. The study archive covered Germany from 2019 to 2024 and provided the observational basis for all models evaluated against individual or spatially aggregated OCO-2 SIF.

\subsection{Sentinel-2 L3A WASP Monthly Composites}

High-resolution optical data were obtained from Sentinel-2 Level 3A monthly composites produced by the Weighted Average Synthesis Processor (WASP). WASP forms temporal syntheses from atmospherically corrected Sentinel-2 Level 2A inputs generated with MAJA, with the processing specifications documented by Hagolle and colleagues \cite{hagolle_maja_2017,hagolle_wasp_2018}. Products were organized by MGRS tile, year, and monthly synthesis date, which provided a consistent spatial reference for the study areas in Germany. The archives supplied blue, green, red, and near-infrared reflectance bands B2, B3, B4, and B8 at 10 m resolution. Red-edge, narrow near-infrared, and short-wave-infrared bands B5, B6, B7, B8A, B11, and B12 were available at 20 m resolution. Each product also contained 10 m and 20 m classification layers, identified as \texttt{FLG\_R1} and \texttt{FLG\_R2}, alongside the reflectance rasters. Reflectance was encoded as signed 16-bit integers with a quantification value of 10,000, and the product documentation designated $-10{,}000$ as the no-data value outside the image. The principal SIF archive contained monthly composites from February to July during 2019--2024, while the later winter-wheat application used available March--July composites for 2019--2022 and 2024.

\subsection{GLASS Vegetation and Radiation Products}

The Global LAnd Surface Satellite (GLASS) product suite supplied medium-resolution vegetation and radiation variables \cite{liang_glass_2013,liang_glass_suite_2020}. The downloaded vegetation products were fraction of absorbed photosynthetically active radiation (FAPAR), normalized difference vegetation index (NDVI), and enhanced vegetation index (EVI). FAPAR was obtained from the GLASS09D01 V60 archive, whereas NDVI and EVI were obtained from the corresponding GLASS13D01 and GLASS14D01 250 m MODIS products. These products were supplied as eight-day composites on the MODIS sinusoidal grid and stored in HDF format. Tiles \texttt{h18v03} and \texttt{h18v04} covered Germany for the spatial domain used in this thesis. The retained archive spanned 2019--2024 and focused on composite start days from day of year 033 through 209, corresponding approximately to February through July. The GLASS photosynthetically active radiation product, GLASS04B01 V42, was additionally acquired at $0.05^{\circ}$ spatial resolution. In contrast to the eight-day vegetation products, the source PAR archive was daily; a reduced eight-day sequence was downloaded for the initial MODIS and polygon-mean experiments.

\subsection{VIIRS Daily PAR Product}

Daily radiation data for the principal Sentinel-2 workflow came from the VIIRS/NPP Photosynthetically Active Radiation Daily Level 3 Global 1 km Sinusoidal Grid Version 2 product, identified as VNP18A2 \cite{wang_vnp18a2_2025}. The product provides daily-average photosynthetically active radiation and an associated categorical quality layer. Its native grid is a VIIRS sinusoidal projection with a pixel spacing of approximately 927 m. The quality layer distinguishes retrievals based on VNP43 surface reflectance, climatological fallback reflectance, non-land surfaces, and missing valid reflectance input. This distinction was important because the quality codes describe data provenance and usability rather than a continuous numerical ranking. Daily-mean PAR was chosen instead of an instantaneous radiation field because the OCO-2 target archive contained daily-corrected SIF. Rasters were acquired for the unique OCO-2 observation dates required by the Sentinel-2 experiments. The local archive stored a daily-mean PAR raster and its corresponding quality raster for each date and year. The source resolution and native coordinate reference system were retained in the downloaded files, while their later alignment with the optical predictor grid is described in Section~\ref{sec:data_preprocessing}.

\subsection{Crop, Regional, and Yield Data}

Annual crop identity was obtained from the German Aerospace Center crop-type map series distributed through the EOC Geoservice \cite{asam_crop_2022,gessner_crop_2025}. The maps provided a separate 10 m categorical raster for each study year and distinguished the main cereal, oilseed, root-crop, grassland, woody-crop, other agricultural, and non-crop classes. These annual rasters supplied the crop information required for footprint composition and the later winter-wheat analysis. Administrative boundaries were obtained from the Geographic Information System of the Commission through the \texttt{giscoR} interface \cite{gisco_nuts}. NUTS level 1 represented German federal states, while NUTS level 3 represented the districts and independent cities used in the Bavarian yield analysis. Sentinel-2 MGRS tile extents provided the separate spatial units used to organize optical raster products. German plant-hardiness-zone polygons supplied the agro-climatic regions used for zone-level model evaluation \cite{germany_hardiness_zones}. Official Bavarian crop-harvest tables provided district-level yields for 2019, 2020, 2021, 2022, and 2024 \cite{bavaria_harvest_statistics}. The selected response was \emph{Winterweizen (einschl. Dinkel und Einkorn)}, which reports winter wheat together with spelt and einkorn. Source values were reported in decitonnes per hectare using German numeric formatting, and 2023 was excluded because matching optical predictor coverage was insufficient.

\section{Data Preprocessing\label{sec:data_preprocessing}}

\subsection{OCO-2 Target and Footprint Preparation}

The OCO-2 files were first converted into a row-wise observation table in which one row represented one sounding. The daily-corrected SIF values at 757 nm and 771 nm were used to construct the main response variable. For sounding $i$, the combined target was defined as

\begin{equation}
    S_i = \frac{S_{757,i}+1.5S_{771,i}}{2},
    \label{eq:combined_sif_target}
\end{equation}

where $S_{757,i}$ and $S_{771,i}$ are the daily-corrected retrievals in the two oxygen absorption bands. The factor of 1.5 brings the 771 nm retrieval closer to the scale of the 757 nm retrieval before the two values are combined. This target is called \texttt{target\_modis\_sif} in the processed tables, although it was also used by Sentinel-2 models. The name was retained for consistency with the earlier workflow and does not mean that the target itself came from MODIS. The 740 nm target was kept only for selected legacy experiments. Using one defined combined target allowed the principal Sentinel-2, current MODIS, and polygon-mean experiments to share the same response variable.

Retrieval uncertainty was adjusted using the OCO-2 daily correction factor before unreliable observations were removed. If $\sigma_{757,i}$ and $\sigma_{771,i}$ denote the corrected uncertainties, the uncertainty of the combined target was calculated as

\begin{equation}
    \sigma_i = \frac{1}{2}
    \sqrt{\sigma_{757,i}^{2}+\left(1.5\sigma_{771,i}\right)^{2}}.
    \label{eq:combined_sif_uncertainty}
\end{equation}

Negative SIF was not automatically treated as invalid because small negative retrievals can occur when the true signal is close to zero. A negative value was accepted when $S_i+2\sigma_i\geq0$, marked as questionable when the two-standard-deviation check failed but $S_i+3\sigma_i\geq0$, and rejected when $S_i+3\sigma_i<0$. The principal datasets retained only observations classified as accepted by this rule. A broad range check of $-0.5\leq S_i<2.0$ was then used to remove clearly implausible combined values. Equivalent checks were stored for the separate 757 nm and 771 nm targets so that they could be selected in earlier multi-target experiments. Quality flag and measurement mode were preserved as separate fields rather than being hidden inside the uncertainty rule. This made it possible to compare nadir and glint observations, quality flags 0 and 1, and different historical subsets later in the analysis.

Each sounding footprint was rebuilt from its four recorded corner coordinates instead of using the central latitude and longitude as its geometry. The corner pairs were ordered to form a closed polygon in WGS 84, and invalid polygons were repaired where possible. Polygon area and centroid operations were carried out after transformation to a projected coordinate system, mainly ETRS89-LAEA Europe (EPSG:3035), so that distances and areas were measured in metres. The centroid was used for unambiguous assignment to states, plant hardiness zones, MGRS tiles, fixed grid cells, and NUTS-3 regions. The polygon itself was retained for raster extraction and for fractional footprint masks. This distinction mattered because an elongated OCO-2 footprint can cross a raster or administrative boundary even when its centroid lies on only one side. Every row also kept a stable row identifier, acquisition date, source file, quality flag, measurement mode, and relevant geographic labels. These fields allowed observations to be traced through later joins, chip manifests, model predictions, and subgroup evaluations.

\subsection{Temporal Matching and Spatial Alignment}

Predictor dates were matched according to the temporal meaning of each raster product. An OCO-2 observation was linked to the Sentinel-2 monthly composite for the same tile, year, and month. The documented lower and upper dates of the Sentinel synthesis were also compared with the OCO-2 date. Observations inside this interval were labelled \texttt{inrange}, while observations outside it were labelled \texttt{outrange}. Both groups were retained in the later main dataset because most out-of-range differences were small and their performance could be reported separately. Eight-day FAPAR, EVI, and NDVI products were matched by the composite interval containing the SIF day rather than by the numerically closest composite start date. For example, a sounding on day of year 007 belongs to the product beginning on day 001, because that composite describes days 001--008. Daily VIIRS PAR was matched to the exact OCO-2 date, while the earlier GLASS PAR branch used the closest available day from the reduced eight-day download sequence.

All image channels within one convolutional sample had to share the same rows, columns, extent, and coordinate reference system. The Sentinel-2 models therefore used a 20 m reference grid in the UTM coordinate system of the relevant Sentinel product. Ten-metre reflectance bands were reduced to this grid by area averaging, while bands already supplied at 20 m remained on their native spacing. Continuous FAPAR and PAR rasters were cropped near the requested chip and reprojected to the same grid using bilinear interpolation. The PAR quality mask was applied on the native VIIRS grid before reprojection, and only quality codes 1 and 2 were accepted. FAPAR values outside the interval from 0 to 1 and PAR values outside the physically screened range from 0 to 700 were treated as missing. APAR was then calculated after alignment as the pixel-wise product of FAPAR and PAR. This ordering ensured that every APAR pixel combined predictor values referring to the same 20 m location.

Categorical information required a different treatment from continuous reflectance and radiation data. The Sentinel classification masks at 10 m and 20 m were used to retain valid land reflectance before the spectral indices were calculated. Annual crop codes were converted into separate binary membership rasters and then averaged onto the 20 m grid. A resulting value of 0.75 therefore means that three quarters of the contributing 10 m area belonged to the stated crop group. No-data crop pixels were kept distinct during source processing, while the final non-crop and valid-data handling followed the requirements of each model-ready export. OCO-2 footprints were rasterized as fractional masks, so boundary pixels could take values between zero and one according to the proportion of the pixel covered by the polygon. Fractional masks were preferred to binary masks because a 20 m or 250 m pixel rarely follows the slanted edge of a sounding exactly. All masks and rasters were clipped to the requested model window after their coordinate systems had been aligned.

\subsection{Predictor and Crop-Channel Construction}

Five spectral indices were calculated from land-masked Sentinel-2 reflectance. Their definitions were

\begin{align}
    \mathrm{NDVI} &= \frac{B8-B4}{B8+B4}, &
    \mathrm{NDMI} &= \frac{B8-B11}{B8+B11}, \nonumber\\
    \mathrm{NDRE} &= \frac{B8A-B5}{B8A+B5}, &
    \mathrm{NIRv} &= B8\times\mathrm{NDVI}, \nonumber\\
    \mathrm{EVI} &= 2.5\frac{B8-B4}{B8+6B4-7.5B2+1}.
    \label{eq:sentinel_indices}
\end{align}

The stored integer reflectance values were divided by 10,000 before these equations were applied, and the designated value of $-10{,}000$ was treated as no-data. A small denominator tolerance prevented division by values close to zero. EVI values outside the interval from $-1$ to 1 were removed because the rational expression can become extremely large when its denominator approaches zero. This check was introduced after a diagnostic found a very small number of extreme EVI pixels that strongly affected normalization. NDVI, NDMI, NDRE, and NIRv retained finite values generated from valid reflectance inputs. These operations produced spatial channels rather than footprint averages for the convolutional models. The same index definitions were also used when polygon means were required for the tabular branch.

The crop information was simplified into channels that represented the main agricultural structures in the study area. The channels were winter wheat, winter barley, winter rye, maize, grass and forage, woody crops, other crops, active crops, and non-crop land. The woody channel combined vineyards with fruit trees and other woody agricultural vegetation, while the grass and forage channel combined the relevant grassland and fodder classes. Active-crop fraction depended on both crop class and observation month. For example, winter wheat was considered active from February to July and again in October and November, while maize was active from May to October. The active channel was formed by summing the fractions of crop classes whose predefined growth period included the current month. This supplied seasonal crop context without asking the network to infer every crop calendar only from the crop label. The annual crop map remained year-specific so that rotations and changes between years were not replaced by one static land-cover map.

Month was represented by two spatially constant channels rather than by the integers 2 through 7. For month $m$, the channels were calculated as

\begin{equation}
    m_{\sin}=\sin\!\left(\frac{2\pi(m-1)}{12}\right),\qquad
    m_{\cos}=\cos\!\left(\frac{2\pi(m-1)}{12}\right).
    \label{eq:cyclical_month}
\end{equation}

This representation places December and January next to each other on a cycle and avoids an artificial large distance between month 12 and month 1. The principal Sentinel-2 chip therefore contained 19 channels: five spectral indices, FAPAR, PAR, APAR, seven grouped crop fractions, active-crop fraction, non-crop fraction, and the two month channels. Latitude and longitude were not included as image channels in this version. Excluding them encouraged the model to use environmental and seasonal information rather than directly memorizing track locations. The channel order was written to the dataset configuration and later checked against the model checkpoint during inference. This was necessary because changing channel order would silently associate trained weights with the wrong predictor.

\subsection{Model-Ready Data Representations}

The processed predictors were organized differently for tabular, area-to-point, and convolutional models. In the polygon-mean branch, every accepted sounding became one tabular row, and the mean of predictor channel $c$ was calculated from the pixels intersecting its footprint:

\begin{equation}
    \bar{x}_{i,c}=
    \frac{\sum_{p\in i}a_{i,p}x_{p,c}}
         {\sum_{p\in i}a_{i,p}},
    \label{eq:polygon_weighted_mean}
\end{equation}

where $a_{i,p}$ is the fractional area of pixel $p$ inside footprint $i$. The resulting table contained one combined SIF target and one value for each selected spectral, radiation, and crop predictor, with APAR approximated as mean FAPAR multiplied by mean PAR. Random Forest and XGBoost used this compact table, which removed within-footprint spatial structure but required relatively little storage. Single-footprint CNN datasets instead centred one multi-channel raster chip on each accepted footprint and stored the footprint as a fractional mask. The predicted SIF map was averaged over this mask before comparison with the observed footprint value. The area-to-point MLP represented each footprint as a set of no more than 200 spatially distributed Sentinel predictor rows and assigned the same footprint-level 740 nm target to those rows. Its pixel predictions were averaged back to the footprint level because directly observed 20 m SIF labels were unavailable. The legacy MLP and initial MODIS 16 by 16 dataset therefore differed from the principal models in both target wavelength and filtering.

Multi-footprint datasets placed several nearby, same-date observations in one predictor window while retaining a separate mask and target for every footprint. The Sentinel-2 version used a 6 km by 6 km window at 20 m resolution, generally contained four to ten footprints, and required at least 90\% of every retained footprint to lie inside the chip. Its errors were averaged within each chip so that chips containing more footprints did not automatically have greater influence. The spatial-aggregate dataset used a different target by assigning same-tile, same-date, and same-mode footprints to fixed 4 km cells containing at least five observations. Each cell had 200 by 200 predictor pixels, and its observed target was

\begin{equation}
    \bar{S}_g=\frac{1}{n_g}\sum_{i=1}^{n_g}S_i,
    \label{eq:spatial_aggregate_target}
\end{equation}

where $n_g$ is the number of assigned footprints and $S_i$ is their observed SIF. The aggregate weight map gave equal influence to the normalized footprint masks even when their areas inside the cell differed. After this model was selected, native-footprint inference chips were prepared for two unseen MGRS tiles using the same 4 km context, 20 m grid, and channel order. The resulting 20 m predictions were averaged over each observed footprint mask, and these external samples were not used to refit the model or its calibration.

For the yield dataset, Bavarian NUTS-3 regions were retained when at least 50\% of their area fell inside an available Sentinel-2 tile, and duplicate tile assignments were removed. Annual crop maps identified winter-wheat pixels for every retained region, year, and month from March to July. A strict 20 m winter-wheat pixel required all four contributing 10 m crop pixels to be classified as winter wheat. CNN inference was limited to blocks containing these pixels, and their predicted SIF values were averaged to one crop-pure monthly value for each NUTS-3 region. Raw comparison variables were formed from all accepted OCO-2 footprints and from subsets with at least 5\%, 10\%, or 30\% winter-wheat coverage. Footprints were first averaged within acquisition date and then across the available dates in each month so that dense OCO-2 tracks did not dominate the regional predictor. The crop-pure and raw monthly values were reshaped into one row per region and year and joined to official winter-wheat yield. Yield was converted from decitonnes per hectare to tonnes per hectare, and rows without a response were removed. Missing predictor values were replaced using means calculated from 2019--2022 only, so the held-out 2024 observations did not affect preprocessing. The completed table allowed crop-pure enhanced SIF and mixed raw OCO-2 SIF to be compared using the same regional yield records.

\section{Model Architectures}

The area-to-point model used a fully connected multilayer perceptron (MLP) that processed one 20 m predictor row at a time. Each row contained 17 features, consisting of Sentinel-2 spectral indices, crop indicators, and two cyclical month variables. The input layer was connected to a hidden layer with 64 neurons, followed by a rectified linear unit (ReLU) activation and dropout with a probability of 0.1. A second hidden layer reduced the representation from 64 to 32 neurons and again applied ReLU and dropout. The final linear layer converted the 32 hidden values into one unconstrained pixel-level SIF estimate. No output activation was used because SIF was treated as a continuous regression target that could include values below zero. The same MLP weights were applied independently to every sampled pixel, so the model did not use the arrangement or neighbourhood of the pixels. The pixel estimates belonging to one OCO-2 footprint were combined through a weighted mean to produce the footprint-level prediction used for supervision.

The single-footprint and multi-footprint convolutional experiments used compact U-Net architectures that preserved the spatial arrangement of the predictor channels. The MODIS models used two encoder levels with 32 and 64 feature maps, followed by a 128-feature bottleneck and a symmetric decoder with 64 and 32 feature maps. Each convolutional block contained two $3\times3$ convolutions, and every convolution was followed by batch normalization and a ReLU activation. The Sentinel-2 multi-footprint model used the same two-level layout but started with 16 feature maps, increased to 32, and used 64 feature maps in the bottleneck. Its convolutional blocks used group normalization and the SiLU activation instead of batch normalization and ReLU. In both versions, $2\times2$ max pooling reduced the spatial dimensions in the encoder, while $2\times2$ transposed convolutions restored them in the decoder. Skip connections concatenated the encoder features with the corresponding decoder features, allowing local boundaries and fine spatial patterns to be recovered after pooling. A final $1\times1$ convolution produced one SIF value for every output pixel without applying a restrictive output activation. The predicted map was then reduced over either one footprint mask or several separate footprint masks, depending on the experiment.

The principal 4 km spatial-aggregate model used a deeper three-level U-Net because its input contained 200 by 200 pixels and therefore supported an additional stage of spatial abstraction. Its 19 predictor channels entered encoder blocks with 16, 32, and 64 feature maps, followed by a bottleneck with 128 feature maps. Every block again used two padded $3\times3$ convolutions with group normalization and SiLU activations. Three max-pooling operations reduced the 200 by 200 input successively to 100 by 100, 50 by 50, and 25 by 25 cells. The decoder reversed this sequence through transposed convolutions and skip connections, using 64, 32, and 16 feature maps before the final $1\times1$ output layer restored a single 200 by 200 SIF map. This architecture allowed the network to combine field-scale texture and crop patterns with broader spatial context while retaining the original 20 m output grid. During model fitting, the output map was multiplied by the normalized aggregate footprint-weight map and summed to one scalar prediction for the 4 km training target. The architecture therefore generated spatially explicit estimates, although the observed supervision remained the aggregated OCO-2 value rather than a directly measured SIF label for every 20 m pixel.

\section{Model Training and Experimental Design}

The data partitions were created before normalization or model fitting, and the split unit was chosen according to the structure of each experiment. For the principal 4 km model, records that shared either an OCO-2 acquisition date or a Sentinel-2 product were connected into one component, and complete components were assigned to training, validation, or testing. This prevented the same acquisition date or predictor image from appearing in more than one partition. Components were allocated to approximate an 80/10/10 split by number of chips while balancing month, measurement mode, and MGRS tile across the three partitions. The Sentinel-2 multi-footprint experiment also used an 80/10/10 split, but its indivisible unit was the complete acquisition date and allocation was performed within month and measurement-mode groups. Unlike the principal 4 km experiment, this date-based split did not additionally prevent the same Sentinel-2 product from occurring in different partitions. The MODIS multi-footprint experiment used a temporal holdout, with 2019--2022 used for training, 2023 for validation, and 2024 for testing, whereas the two MODIS single-footprint experiments used a 70/15/15 stratified random split by state, year, and MODIS composite day. All sampled pixels from one footprint remained together in the area-to-point MLP split, which used 70\% for training, 15\% for validation, and 15\% for testing. The polygon-mean tree models used a date-grouped 80/20 train/test split stratified by month, measurement mode, and target-value class, while the yield models used 2019--2022 region-year records for training and retained all 2024 records as an independent temporal test set.

Normalization parameters were calculated from the training partition only and then applied unchanged to validation and test data. A random seed of 42 was used for the stochastic data partitions and model training procedures. Raster channels were standardized separately; missing raster values were replaced by zero only after standardization, so that zero represented the training-channel mean rather than a physical zero. Fractional footprint and aggregate masks were converted to floating-point values and normalized to sum to one before they were used to reduce a predicted SIF map to its supervised scalar value. The SIF targets used by the CNN models were also standardized during optimization and returned to their original units for evaluation, whereas the area-to-point MLP retained its target in the original SIF units. The convolutional models and the MLP were optimized with AdamW using an initial learning rate of $10^{-3}$ and a weight decay of $10^{-4}$. The two Sentinel-2 U-Nets used batches of eight chips and accumulated gradients over four batches, giving an effective batch size of 32 without requiring all 32 large images to be held in GPU memory simultaneously. When a chip contained several footprints, the footprint losses were first averaged within that chip and the chip losses were then averaged across the batch, which prevented chips containing more footprints from receiving greater influence solely because of their footprint count. Huber loss was used for the corrected combined-SIF neural experiments, while the original 740 nm baseline used mean squared error; the area-to-point MLP used Huber loss with a delta of 0.1. The Random Forest and XGBoost models were trained with fixed settings so that the comparison focused on the input representation and model family rather than an unequal amount of hyperparameter search.

Neural-network checkpoints were selected using the lowest validation RMSE. The Sentinel-2 models used early stopping with a patience of seven epochs and reduced the learning rate when validation RMSE stopped improving, while the area-to-point MLP allowed up to 150 epochs with a patience of 20. The MODIS networks used a fixed number of epochs but retained the checkpoint with the lowest validation RMSE. For the corrected combined-SIF CNN experiments, a linear calibration of observed SIF against predicted SIF was fitted on the validation predictions only and was then applied unchanged to the test predictions. Raw and calibrated metrics were retained so that any improvement caused by this post-hoc adjustment remained visible. The external native-footprint experiment did not train another network; it applied the saved 4 km U-Net, its training normalization statistics, and its validation-derived calibration to centred 4 km windows from two previously unseen MGRS tiles. The resulting 20 m prediction maps were averaged over the individual OCO-2 footprint masks and compared with their observed footprint-mean SIF values. The winter-wheat yield experiment used the same fixed Random Forest configuration for all crop-pure and raw-SIF predictor variants, without a validation-based tuning stage, so that the 2024 comparison reflected the change in SIF representation. Table~\ref{tab:model-training-settings} summarizes the settings used in the final experiments.

\begin{table}[p]
\centering
\caption[Training and partition settings]{Training and partition settings used in the final modelling experiments. \textit{Notes:} Component splitting kept observations sharing an OCO-2 date or Sentinel-2 product in the same partition. Date and footprint splitting kept complete acquisition dates and footprints together, while stratified splitting used state, year, and composite day. Year entries show training/validation/test periods, except for the yield model, which used a training/test split. Batch $8\times4$ denotes a batch size of eight with four-step gradient accumulation, and ES denotes early-stopping patience. All neural models used a weight decay of $10^{-4}$ and retained the checkpoint with the lowest validation RMSE. External native-footprint inference is omitted because it reused the trained 4 km checkpoint without further fitting.}
\label{tab:model-training-settings}
\small
\renewcommand{\arraystretch}{1.15}
\resizebox{\textwidth}{!}{%
\begin{tabular}{l l l l c c c c c}
\hline
\textbf{Experiment} & \textbf{Split} & \textbf{Loss} & \textbf{Optimizer} &
\textbf{Batch} & \textbf{Epochs/trees} & \textbf{LR} &
\textbf{ES} & \textbf{Calibration} \\
\hline
4 km aggregate U-Net & Component 80/10/10 & Huber (1.0) & AdamW & $8\times4$ & 30 & $10^{-3}$ & 7 & Yes \\
Sentinel-2 multi-SIF U-Net & Date 80/10/10 & Huber (1.0) & AdamW & $8\times4$ & 30 & $10^{-3}$ & 7 & Yes \\
MODIS multi-SIF U-Net & Year 2019--22/23/24 & Huber (0.5) & AdamW & 128 & 15 & $10^{-3}$ & -- & Yes \\
MODIS single-SIF U-Net & Stratified 70/15/15 & Huber (0.5) & AdamW & 256 & 30 & $10^{-3}$ & -- & Yes \\
MODIS 740 nm baseline & Stratified 70/15/15 & MSE & AdamW & 512 & 30 & $10^{-3}$ & -- & No \\
Area-to-point MLP & Footprint 70/15/15 & Huber (0.1) & AdamW & -- & 150 & $10^{-3}$ & 20 & No \\
Polygon means: XGBoost & Date 80/20 & Squared error & XGBoost & -- & 800 & 0.03 & -- & No \\
Polygon means: Random Forest & Date 80/20 & Squared error & RF & -- & 700 & -- & -- & No \\
Winter-wheat yield RF & Year 2019--22/24 & Squared error & RF & -- & 500 & -- & -- & No \\
\hline
\end{tabular}
}
\end{table}

\section{Experimental Environment and Reproducibility}

All data acquisition, preprocessing, spatial matching, raster extraction, and model-ready dataset construction were performed on a local workstation with a 20-core processor and 32 GB of memory. Python processing used Python 3.11.9 in the Positron development environment, with NumPy and pandas for tabular operations and Rasterio, GeoPandas, Shapely, and pyproj for raster, vector, geometry, and coordinate-reference-system operations. R processing used R 4.5.2 in RStudio, with \texttt{terra} and \texttt{sf} providing the main raster and vector functionality and the \texttt{future}, \texttt{furrr}, and tidyverse packages supporting parallel and tabular workflows. Computationally expensive operations were reduced by grouping observations by date, product, or tile, cropping rasters before extraction, and processing independent groups in parallel where possible. Neural-network training and evaluation were conducted in Kaggle environments providing two NVIDIA T4 GPUs, a four-core Intel Xeon CPU allocation, and 30 GB of memory. The Kaggle models used PyTorch 2.10.0 with CUDA 12.8, although the availability of two GPUs does not imply that every experiment used distributed multi-GPU training. Random seeds, data-partition definitions, normalization statistics, chip metadata, prediction tables, and evaluation outputs were stored with the corresponding scripts and notebooks so that each experiment could be reconstructed from its model-ready inputs. The final retained model artifact was the checkpoint and associated weights for the best-performing 4 km spatial-aggregate U-Net, together with the normalization parameters and channel order required to apply it to new data.

\section{Techniques for Visualization}

\section{Evaluation Metrics}

Let $y_i$ denote the observed SIF or yield value for test sample $i$, let $\widehat{y}_i$ denote its prediction, and define the residual as $e_i=\widehat{y}_i-y_i$. Predictive error was summarized using root mean squared error (RMSE), mean absolute error (MAE), and mean bias error, while the coefficient of determination ($R^2$) measured the fraction of observed variation explained by the predictions. A positive bias indicates systematic overprediction under the adopted residual convention, whereas a negative value indicates underprediction. Lower RMSE and MAE values indicate predictions that are closer to the observations, with RMSE placing greater emphasis on large errors. An $R^2$ value near one indicates strong agreement in variation, while a value below zero means that the predictions perform worse than using the observed mean as a constant estimate. Because RMSE depends on the spread of the target, normalized RMSE (NRMSE) was also used when comparing targets defined at different spatial supports and was calculated by dividing RMSE by the standard deviation of the observed test values, $s_y$. These metrics were calculated as
\begin{equation}
\begin{aligned}
    \mathrm{RMSE} &= \sqrt{\frac{1}{n}\sum_{i=1}^{n}\left(\widehat{y}_i-y_i\right)^2}, \\
    \mathrm{MAE} &= \frac{1}{n}\sum_{i=1}^{n}\left|\widehat{y}_i-y_i\right|, \\
    \mathrm{Bias} &= \frac{1}{n}\sum_{i=1}^{n}\left(\widehat{y}_i-y_i\right), \\
    R^2 &= 1-\frac{\sum_{i=1}^{n}\left(y_i-\widehat{y}_i\right)^2}
    {\sum_{i=1}^{n}\left(y_i-\overline{y}\right)^2}, \\
    \mathrm{NRMSE} &= \frac{\mathrm{RMSE}}{s_y}.
\end{aligned}
\label{eq:evaluation_metrics}
\end{equation}

Huber loss was used as the optimization objective for the corrected combined-SIF neural-network experiments and the area-to-point MLP, with the transition parameter varying by experiment as summarized in Table~\ref{tab:model-training-settings}. It behaves quadratically for small residuals but changes to a linear penalty for large residuals, making model fitting less sensitive to unusually large SIF errors than mean squared error. For a residual $e_i=\widehat{y}_i-y_i$ and transition parameter $\delta$, the loss was defined separately as
\begin{equation}
    L_{\delta}(e_i)=
    \begin{cases}
        \dfrac{1}{2}e_i^2, & \text{if } |e_i|\leq\delta, \\
        \delta\left(|e_i|-\dfrac{1}{2}\delta\right), & \text{if } |e_i|>\delta.
    \end{cases}
    \label{eq:huber_loss}
\end{equation}
Huber loss was used to update model parameters during training, whereas RMSE, MAE, bias, and $R^2$ remained the principal metrics used to report predictive performance on validation and test data.

Agreement was examined further using the slope and intercept of the least-squares regression of predicted values on observations, Pearson's correlation coefficient, and the distribution of residuals. The regression slope $b_1$ and Pearson correlation $r$ were calculated as
\begin{equation}
\begin{aligned}
    b_1 &=
    \frac{\sum_{i=1}^{n}(y_i-\overline{y})(\widehat{y}_i-\overline{\widehat{y}})}
    {\sum_{i=1}^{n}(y_i-\overline{y})^2}, \\
    b_0 &= \overline{\widehat{y}}-b_1\overline{y}, \\
    r &=
    \frac{\sum_{i=1}^{n}(y_i-\overline{y})(\widehat{y}_i-\overline{\widehat{y}})}
    {\sqrt{\sum_{i=1}^{n}(y_i-\overline{y})^2
    \sum_{i=1}^{n}(\widehat{y}_i-\overline{\widehat{y}})^2}}.
\end{aligned}
\label{eq:agreement_metrics}
\end{equation}
An ideal prediction has $b_0=0$, $b_1=1$, zero bias, and small residual spread; residual stability was therefore also described by its standard deviation and interquartile range, $Q_{0.75}(e)-Q_{0.25}(e)$. Predictor importance for the principal U-Net was estimated by independently permuting each validation channel and recording the increase in RMSE, $I_j=\mathrm{RMSE}_{\mathrm{perm}(j)}-\mathrm{RMSE}_{\mathrm{base}}$, with larger positive values indicating greater predictive importance. The same core metrics were recalculated for SIF-value bins, footprint counts, measurement modes, months, years, MGRS tiles, states, acquisition-date alignment classes, and plant-hardiness zones to identify systematic differences among data groups. For the winter-wheat comparison, the relative RMSE reduction of crop-pure SIF against a raw-SIF baseline was calculated as $100(\mathrm{RMSE}_{\mathrm{raw}}-\mathrm{RMSE}_{\mathrm{crop}})/\mathrm{RMSE}_{\mathrm{raw}}$. Paired bootstrap resampling by NUTS-3 region preserved the correspondence between competing yield-model predictions, and the 2.5th and 97.5th percentiles of the bootstrap improvement distribution formed the reported 95\% confidence interval. Spatial residual maps and observed-versus-predicted plots complemented the numerical metrics by showing where errors occurred and whether the prediction range was compressed relative to the observations.

%---------------------------------
